% part: intuitionistic-logic
% chapter: propositions-as-types
% section: sequent-natural-deduction

\documentclass[../../../include/open-logic-section]{subfiles}

\begin{document}

\olfileid{int}{pty}{snd}

\olsection{Dedhuksi Alami Sekuen}

Kita nulis $\Gamma \Sequent !A$ yen ana !!{derivation} dedhuksi alami
kanthi $\Gamma$ minangka asumsi !!{undischarged} lan $!A$ minangka
konklusi; utawa $\Sequent !A$ yen $\Gamma$ kosong.

Kita nulis $\Gamma, !A_1, \dots, !A_n$ kanggo $\Gamma \cup \{!A_1,
\dots, !A_n\}$, lan $\Gamma, \Delta$ kanggo $\Gamma \cup \Delta$.

Gatekna yen kita nduweni $\Gamma \Sequent !A \land !B$, tegese ana
!!a{derivation} kanthi $\Gamma$ minangka asumsi !!{undischarged} lan
$!A \land !B$ minangka !!{formula} pungkasan. Kanthi ngetrapake
$\Elim{\land}$ ing ngisor, kita bisa ngasilake !!a{derivation} kanthi
asumsi !!{undischarged} sing padha lan $!A$ minangka konklusi.
Kanthi tembung liya, yen $\Gamma \Sequent !A \land !B$, mula
$\Gamma \Sequent !A$.
\begin{prooftree}
  \Axiom$\Gamma \fCenter !A \land !B$
  \RightLabel{$\Elim{\land}$}
  \UnaryInf$\Gamma \fCenter !A$
  \DisplayProof\qquad\bottomAlignProof
  \Axiom$\Gamma \fCenter !A \land !B$
  \RightLabel{$\Elim{\land}$}
  \UnaryInf$\Gamma \fCenter !B$
\end{prooftree}
Label $\Elim{\land}$ nuduhake sesambungane karo aturan kanthi jeneng
sing padha ing dedhuksi alami.

Semono uga, upamane kita nduweni $\Gamma, !A \Sequent !B$, tegese ana
!!a{derivation} kanthi asumsi !!{undischarged} $\Gamma, !A$ lan
!!{formula} pungkasan~$!B$. Yen ngetrapake aturan $\Intro{\lif}$, kita
nduweni !!a{derivation} kanthi $\Gamma$ minangka asumsi !!{undischarged}
lan $!A \lif !B$ minangka !!{formula} pungkasan, yaiku
$\Gamma \Sequent !A \lif !B$. Gatekna yen iki nggawe !!{discharge}
asumsi luwih cetha.
\begin{prooftree}
  \Axiom $\Gamma, !A \fCenter !B$
  \RightLabel{$\Intro{\lif}$}
  \UnaryInf $\Gamma \fCenter !A \lif !B$
\end{prooftree}

Kita bisa narik konklusi saka aturan liyane kanthi cara sing padha,
sing diandharake kaya mangkene:
\begin{gather*}
  \Axiom $\Gamma \fCenter !A$
  \Axiom $\Delta \fCenter !B$
  \RightLabel{$\Intro{\land}$}
  \BinaryInf $\Gamma,\Delta \fCenter !A \land !B$
  \DisplayProof\\
  \Axiom $\Gamma \fCenter !A \land !B$
  \RightLabel{$\Elim{\land}_1$}
  \UnaryInf $\Gamma \fCenter !A$
  \DisplayProof
  \qquad
  \Axiom $\Gamma \fCenter !A \land !B$
  \RightLabel{$\Elim{\land}_2$}
  \UnaryInf $\Gamma \fCenter !B$
  \DisplayProof
  \\
  \Axiom $\Gamma \fCenter !A$
  \RightLabel{$\Intro{\lor}_1$}
  \UnaryInf $\Gamma \fCenter !A \lor !B$
  \DisplayProof\qquad
  \Axiom $\Gamma \fCenter !B$
  \RightLabel{$\Intro{\lor}_2$}
  \UnaryInf $\Gamma \fCenter !A \lor !B$
  \DisplayProof\\
  \Axiom $\Gamma \fCenter !A \lor !B$
  \Axiom $\Delta, !A \fCenter !C$
  \Axiom $\Delta', !B \fCenter !C$
  \RightLabel{$\Elim{\lor}$}
  \TrinaryInf $\Gamma, \Delta, \Delta' \fCenter !C$
  \DisplayProof
  \\
  \Axiom $\Gamma, !A \fCenter !B$
  \RightLabel{$\Intro{\lif}$}
  \UnaryInf $\Gamma \fCenter !A \lif !B$
  \DisplayProof
  \qquad
  \Axiom $\Delta \fCenter !A \lif !B$
  \Axiom $\Gamma \fCenter !A$
  \RightLabel{$\Elim{\lif}$}
  \BinaryInf $\Gamma, \Delta \fCenter !B$
  \DisplayProof
  \\
  \Axiom $\Gamma \fCenter \lfalse$
  \RightLabel{$\FalseInt$}
  \UnaryInf $\Gamma \fCenter !A$
  \DisplayProof
\end{gather*}

Asumsi dhewe iku !!a{derivation} saka $!A$ adhedhasar~$!A$, yaiku
kita tansah nduweni $!A \Sequent !A$.

\begin{prooftree}
  \AxiomC{$ $}
  \UnaryInfC{$!A \fCenter !A$}
\end{prooftree}

Aturan iki bebarengan bisa dianggep kalkulus ngenani !!{derivation}s
dedhuksi alami sing ana. Aturan iki uga bisa dianggep varian notasi
dedhuksi alami, kanthi saben langkah nyathet ora mung !!{formula}
sing !!{derive}, nanging uga asumsi !!{undischarged} sing dadi
dhasar !!{derive} kasebut.

\begin{prooftree}
  \Axiom$!A \fCenter !A$
  \UnaryInf $!A \fCenter !A \lor (!A \lif \lfalse)$
  \Axiom $!B \fCenter !B$
  \BinaryInf$!A, !B \fCenter \lfalse$
  \UnaryInf$!B \fCenter !A \lif \lfalse$
  \UnaryInf$!B \fCenter !A \lor (!A \lif \lfalse)$
  \Axiom$!B \fCenter !B$
  \BinaryInf$!B \fCenter \lfalse$
  \UnaryInf$\fCenter !B \lif \lfalse$
\end{prooftree}
kanthi $!B$ minangka cekakan saka $(!A \lor (!A \lif \lfalse))\lif \lfalse$.

\end{document}
